\section{Outcome, scope, and the audited baseline}\label{sec:scope}

The immediate outcome is a new family of exact Jones obstruction backends,
together with a component factorization that improves their dependence on
the geometry of a supplied crossing order. These are useful pieces of a
complete recognizer. They are not, by themselves, complete recognition
algorithms: inequality with the unknot's Jones value proves knottedness,
whereas equality requires another method.

This distinction determines both the software interface and the mathematical
claims. In the implementation, a successful obstruction returns
\code{KNOTTED}. Equality returns an inconclusive result. Exhausting a local
filter budget also returns an inconclusive result, with its reason recorded;
it does not produce a topological verdict. The enclosing recognizer retains
its existing complete Khovanov fallback, subject to the user's global resource
limits. Thus a bounded execution can still return \code{UNKNOWN}.

\subsection{What was inspected and preserved}

The baseline is ProveIt commit
\code{86f23490006c880cd01db9bec10761d39c6395c8}, dated
8 October 2026 \cite{proveit2026}. The source snapshot includes the Python
\code{fastunknot} package, its fixtures and tests, the synthesis sources, and
the report summaries. Downloaded source blobs were checked against their
Git object identifiers. The included reference implementation is immutable
within the experiment package. The baseline test suite passed all 134 tests.

The baseline already contains substantial work: short-braid certificates,
Seifert and descending-diagram tests, Reidemeister simplification,
connected-sum handling, modular Alexander and Jones obstructions,
optional exact Alexander computation, and several exact Khovanov scanner
representations. Earlier research also introduced canonical multiplicities,
shared morphisms, radical perturbation inverses, saturated or Euler-assisted
terminal decisions, homological windows, and scalar Fitting decompositions
\cite{report07,report08,report09,report10,report11,report12}.
Those ideas remain predecessors, rather than contributions claimed here.
The immediately preceding continuation's scalar splitting and
common-nilpotent decision shortening were inspected separately; the
ring-valued extension in Section~\ref{sec:residue} is carefully distinguished
from them.

The current Jones filter stores boundary smoothing matchings and their
scalar coefficients. For an arbitrary cut with $w$ endpoints, a universal
matching count is $(w-1)!!$, not a Catalan number. A Catalan bound requires
additional boundary-order geometry that an arbitrary crossing prefix need
not possess. This creates an opportunity to change the scalar representation
without making an invalid noncrossing assumption.

\subsection{Claims and their evidential status}

\begingroup\small
\begin{longtable}{@{}>{\raggedright\arraybackslash}p{0.28\textwidth}p{0.67\textwidth}@{}}
\toprule
Claim & Status and limit\\
\midrule
\endfirsthead
\toprule Claim & Status and limit\\ \midrule
\endhead
Signed Tait/Potts expression for Jones values &
Classical; derived in the repository's normalization and checked against an
independent smoothing cube.\\
At most $w/2$ active Tait vertices &
Proved here by a direct boundary-transition count for every crossing order
of a connected plane knot shadow.\\
Fixed-$q$ equality-pattern transfer &
Implemented and verified; costs $\poly(n,w)q^{w/2}$, with linear-in-$n$
coefficient bit lengths for fixed $q$.\\
Independent-component transfer &
Implemented; a relative-color merge restores precisely the information
needed when components first interact.\\
$q=5$ blind weaving family &
Proved by specializing the classical three-braid trace formula; preserved as
an adversarial regression family.\\
$q\ge6$ detects that family &
Proved; says nothing universal about other knot families.\\
Residue-polynomial Fitting certificates &
Proved for matrices over specified commutative truncated rings; verified by
a separate algebraic oracle.\\
Quasi-polynomial full recognition &
Conditional research goal. No complete complexity proof or complete
quasi-polynomial implementation is supplied.\\
\bottomrule
\caption{The claim ledger. A proof supplied in this article is not a claim
of literature-wide priority.}
\label{tab:claims}\\
\end{longtable}
\endgroup

\subsection{Prior literature and the nature of the improvement}

The Potts formulation of the Jones polynomial belongs to the classical
state-model literature \cite{kauffman1987,sokal2005,moffatt2011}.
Meichanetzidis and Kourtis explicitly used the same integer-color
specializations, including $q=5$, in tensor-network evaluation
\cite{meichanetzidis2019}. Makowsky established bounded-treewidth
tractability for signed Tutte and Kauffman computations
\cite{makowsky2005}. Maria gave single-exponential width dependence for
general quantum link invariants \cite{maria2021}. The broad statement
that Jones evaluation is single exponential in a suitable width is therefore
not a new result.

The contribution appropriate to this continuation is more concrete. We
derive a bound directly in the crossing-order parameter used by ProveIt,
make equality-pattern aggregation and normalization explicit, implement exact
quadratic arithmetic, retain resource-safe recognition semantics, identify a
structural failure of the $q=5$ point, and supply measurements including
regressions. The component factorization and residue-algebra sections give
additional, independently checkable improvements whose exact applicability
is stated.

The public \code{openai/math} collection was inspected as suggested by the
user \cite{openaimath2026}. It did not supply a missing unknot-recognition
theorem. Its useful role here is methodological: formulate explicit lemmas,
test a proposed general claim on adversarial families, and preserve enough
artifacts that independent readers can reproduce the calculation. None of
the results in this article is attributed to that collection without a
specific mathematical dependency.

\subsection{Complexity terminology}

Throughout, $n$ is the number of crossings in the supplied diagram, $w$ is a
maximum cut-edge count for a crossing order, and $q$ is an integer color
count. A statement for fixed $q$ does not hide a growing input parameter.
We use \emph{quasi-polynomial} in the broad sense
$\exp((\log(n+2))^{O(1)})$, and state the stronger bound
$n^{O(\log n)}$ explicitly when it is obtained. Arithmetic-operation bounds
and bit-operation bounds are separated. Timing improvements on a finite
corpus are not asymptotic evidence by themselves.
